% English Abstract

With the rapid development of 5G communications and streaming technologies, video data has become the dominant component of global Internet traffic. Achieving high-quality visual experiences under limited bandwidth has become a critical challenge. However, traditional video coding standards face a trade-off between compression efficiency and computational complexity, while end-to-end deep learning compression schemes, despite their excellent performance, are difficult to deploy on edge devices.

To address these challenges, this thesis proposes a reinforcement learning-based co-processing solution for video coding. The proposed system is built around the NVIDIA Video Encoder (NVENC) hardware encoder, incorporating a lightweight Rate-Adaptive Pre-processing Network (RAPNet) to perform intelligent video enhancement before encoding. The system employs an Actor-Critic architecture with a residual learning strategy that constrains the action space to fine-grained corrections of the original frames. A compound reward function based on relative rate-distortion gain is designed to achieve effective trade-offs between perceptual quality and transmission bitrate. Furthermore, a composite loss strategy comprising policy gradient loss, pixel-level consistency loss, and total variation smoothness loss is proposed to effectively suppress high-frequency artifacts during the enhancement process.

Experiments are conducted on 1080p resolution video sequences in YUV420 format at multiple bitrate points including 800Kbps, 400Kbps, and 200Kbps. Results demonstrate that the proposed scheme achieves approximately 5\% BD-Rate optimization in LPIPS and SSIM metrics compared to the baseline, with the model parameter count strictly controlled under 1M, effectively improving the perceptual quality of videos at low bitrates.

% Note: English keywords are set in \whusetup within my_paper.tex